
\documentclass[UTF8]{ctexart}

\usepackage{amsmath,amssymb,bm}
\usepackage{geometry}
\geometry{a4paper,margin=2.5cm}

\title{动态禁飞区时空解析验证模型}
\author{}
\date{}

\begin{document}
\maketitle

\section{建模思想}

对于任意两个巡检任务节点 $i,j$，考虑无人机从节点 $i$ 到节点 $j$ 的飞行航段。
动态禁飞区同时具有空间约束和时间约束，因此产生冲突需要同时满足：

\begin{enumerate}
    \item 无人机轨迹进入禁飞区域；
    \item 进入禁飞区域的时间位于禁飞时间窗口内。
\end{enumerate}

因此建立空间--时间联合约束模型。

\section{航段轨迹参数化}

设无人机由节点
\[
A=(x_i,y_i)
\]
飞往
\[
B=(x_j,y_j).
\]

航段长度：

\[
L_{ij}=\sqrt{(x_j-x_i)^2+(y_j-y_i)^2}.
\]

若采用原始坐标单位，则无人机速度为：

\[
v=55\mathrm{km/h}=550\text{坐标单位}/h.
\]

因此飞行时间：

\[
\tau_{ij}=\frac{L_{ij}}{550}.
\]

令无人机出发时刻为 $t_i$，定义轨迹参数：

\[
\bm p(\lambda)=
\bm A+\lambda(\bm B-\bm A),
\qquad 0\leq\lambda\leq1 .
\]

对应时间：

\[
t=t_i+\lambda\tau_{ij}.
\]

该参数化保证无人机沿直线匀速飞行。

\section{空间冲突解析判断}

设圆形禁飞区：

\[
\Omega_q=
\{(x,y):\|\bm p-\bm c_q\|\leq R_q\}.
\]

其中圆心：

\[
\bm c_q=(x_q,y_q).
\]

判断航段与禁飞区边界：

\[
\|\bm p(\lambda)-\bm c_q\|^2=R_q^2.
\]

展开得到二次方程：

\[
a\lambda^2+b\lambda+c=0,
\]

其中：

\[
a=(\bm B-\bm A)\cdot(\bm B-\bm A),
\]

\[
b=2(\bm A-\bm c_q)\cdot(\bm B-\bm A),
\]

\[
c=(\bm A-\bm c_q)\cdot(\bm A-\bm c_q)-R_q^2.
\]

求解得到交点参数
$\lambda_1,\lambda_2$。

若：

\[
[\lambda_1,\lambda_2]\cap[0,1]\neq\varnothing,
\]

则该航段在空间上进入禁飞区域。

\section{时间冲突判断}

由空间交点得到进入和离开禁飞区的时间：

\[
t_1=t_i+\lambda_1\tau_{ij},
\]

\[
t_2=t_i+\lambda_2\tau_{ij}.
\]

设禁飞时间窗口：

\[
T_q=[T_q^{s},T_q^{e}],
\]

则动态冲突条件为：

\[
[t_1,t_2]\cap[T_q^s,T_q^e]\neq\varnothing.
\]

即：

\[
\boxed{
\text{空间相交}
+
\text{时间相交}
=
\text{禁飞冲突}
}
\]

\section{动态航段代价}

定义：

\[
\widetilde{\tau}_{ij}(t)
\]

表示无人机在时刻 $t$ 从节点 $i$ 到节点 $j$
的最短合法飞行时间。

第一、二问：

\[
\tau_{ij}=\frac{d_{ij}}{55}.
\]

第三问：

\[
\widetilde{\tau}_{ij}(t)
=
\min
\{
\tau_{\mathrm{direct}},
\tau_{\mathrm{wait}},
\tau_{\mathrm{detour}}
\}.
\]

其中分别对应：

\begin{itemize}
    \item 直接飞行；
    \item 等待禁飞区结束；
    \item 绕行禁飞区。
\end{itemize}

通过该函数可将动态禁飞约束嵌入多无人机路径优化模型。

\section{模型优势}

相比时间离散采样方法，本模型：

\begin{enumerate}
    \item 不存在时间步长导致的漏检问题；
    \item 通过解析求交获得精确冲突时间；
    \item 每条航段仅需常数次计算；
    \item 适合嵌入启发式或精确优化算法。
\end{enumerate}

\end{document}
